\chapter*{General Introduction}
\phantomsection\addcontentsline{toc}{chapter}{General Introduction}
\label{ch:context}
\begingroup
\clubpenalty=10000
\introsection{Overview}
Agile quadrotors are increasingly used as mobile Internet-of-Things (IoT) platforms for inspection, monitoring, mapping, and data collection. However, their lightweight structure and operation in open environments make them particularly vulnerable to wind disturbances, which can degrade navigation performance and increase collision risk. Existing work often addresses either disturbance rejection through adaptive control methods such as $\mathcal L_1$ and incremental nonlinear dynamic inversion (INDI), or safety through control barrier functions (CBFs) and data-driven safe reinforcement learning, while comparatively few approaches treat robustness and safety jointly. This report examines how these directions can be combined in a navigation architecture: reinforcement learning supplies task-oriented decisions, while a control-theoretic safety layer restricts the commands applied to the vehicle. The discussion begins with the basic language of dynamical systems, introduces reinforcement learning and constrained policy optimization, and then develops the role of control barrier functions. It closes by reviewing approaches to disturbance rejection and online adaptation. The central motivation is to understand what each component contributes to reliable navigation when the operating conditions differ from those used for training.

Aerial mobility is particularly useful when sensors must reach locations that are difficult to instrument permanently. An inspection vehicle can approach different parts of a structure, while a monitoring mission can revisit selected locations as conditions change. In both cases, mobility links the physical task to information collection: the robot must reach a useful observation position and remain sufficiently well controlled to collect meaningful measurements. Reliability therefore concerns the continuity and quality of the mission as well as the vehicle's motion~\cite{mozaffari2019tutorial,bouhamed2020ddpg}.

\introsection{Background and Motivation}
Agile quadrotors are increasingly deployed as mobile Internet-of-Things (IoT) platforms for infrastructure inspection, environmental sensing, mapping, and data collection~\cite{mozaffari2019tutorial,bouhamed2020ddpg}. In these applications, mission success depends not only on reaching a destination but also on preserving reliable sensing and data acquisition in uncertain environments. A navigation error near an obstacle may interrupt the mission, damage the vehicle or its payload, and leave the collected data incomplete.

Modern quadrotor navigation stacks often combine map-based planning with a low-latency local controller or policy. Mapping and hierarchical planning maintain feasible routes as the environment is explored~\cite{zhou2021fuel}, whereas learning-based navigation can map compact observations directly to control commands~\cite{bouhamed2020ddpg,kaufmann2023champion}. These components address complementary needs, but they also have different failure modes: a planned route may become stale between updates, while a learned policy may extrapolate poorly beyond the conditions represented during training.

Outdoor operation introduces an additional mismatch between commanded and realized motion. Wind, aerodynamic effects, model mismatch, and low-level tracking errors perturb the translational response and can sharply reduce obstacle clearance~\cite{faessler2018drag}. Learned residual models have enabled agile flight in strong winds~\cite{litNeuralFly2022}, while disturbance-observer methods reconstruct external effects from measured motion and incorporate them into robust safety constraints~\cite{wang2023dobcbf,kalaria2024dobcbfrl}. Such methods motivate adaptation during operation, but their assumptions, data requirements, or preselected robustness levels may limit portability across changing conditions.

\introsection{Related Work}
Several learning and control approaches address navigation safety. Constrained policy optimization can include safety requirements during training, although it generally requires substantial interaction data~\cite{litSafePrimalDual2023}. At deployment, control barrier function (CBF) filters evaluate each proposed command against the current safety condition and modify it only when necessary~\cite{litAmes2017,litAmes2019}. CBF-RL exposes the policy to this correction during training through action substitution without changing the RL optimizer~\cite{litCBFRL2026}. Complementary risk-aware planning and filtering methods use conditional value-at-risk (CVaR) to emphasize rare but consequential tail events~\cite{rockafellar2000cvar,hakobyan2022wasserstein,safaoui2024drcvar,chriat2024wcvarcbf}. Wasserstein distributionally robust optimization (DRO) extends this perspective by guarding against mismatch between empirical and operating distributions~\cite{thEsfahani2018,long2024sensor,mandaokar2026druav}. 

A remaining challenge is the gap between training-time disturbance assumptions and deployment-time conditions. A fixed safety margin can restrict useful movement under mild conditions yet leave insufficient clearance when the disturbance becomes stronger. Likewise, a policy that performs well across a training distribution may encounter persistent wind or combinations of conditions that were poorly represented during learning. This motivates examining both the disturbance response of the controller and the assumptions behind its safety mechanism.

The report follows three connected questions. First, how can a policy learn useful navigation behavior from interaction? Second, how can a separate control mechanism restrict commands before they are applied? Third, how can the overall system respond when its nominal motion model becomes inaccurate? These questions distinguish task performance, constraint enforcement, and adaptation while keeping their interaction visible.

The distinction matters in an inspection mission near a wall. A route planner may provide a collision-free reference, yet wind can displace the vehicle from that route. A tracking controller can reduce the displacement, but its transient response may still consume part of the available clearance. A safety filter must therefore reason about the motion that can actually occur, rather than considering the planned path alone. This example motivates the progression from feedback control to learned policies, barrier filtering, and disturbance-aware navigation.

\introsection{Contribution}
Navigation in changing wind calls for a safety margin that follows the robot's actual motion. Long et al. account for uncertainty using updated sensor measurements and state estimates~\cite{long2024sensor}. Their sensor-based construction does not explicitly track a rolling history of disturbance prediction errors. This leaves the adaptation to changing motion mismatch to a separate mechanism. We address this gap by making recent disturbance residuals directly determine the CBF safety margin.

Our approach combines an EMA-inspired disturbance observer with a rolling residual window. The observer captures the recent disturbance trend; the window retains the errors that remain after estimation. At each update, we rebuild their empirical distribution and compute a Wasserstein--CVaR allowance in the direction of the current obstacle. The safety layer therefore responds to the mismatch observed during operation, while the navigation policy remains fixed.

This design uses measurements already available in the control loop and maintains a bounded history. It requires neither policy retraining nor a separately fitted residual dynamics network. The distributional calculation produces a scalar margin before the command projection, keeping the online adaptation compact. For connected robots with shared sensing and computing resources, this is a direct way to incorporate changing disturbances into action filtering while keeping the allowance tied to recent prediction errors.

\introsection{Report Organization}
Chapter~\ref{ch:stateart} presents the state of the art in feedback control, reinforcement learning, barrier-based safety and disturbance rejection. Chapter~\ref{ch:system} defines the system model and navigation problem. Chapter~\ref{ch:contribution} develops the adaptive Wasserstein--CVaR barrier filter. Chapter~\ref{chap:results} presents the simulation setup, evaluation results and practical sim2sim extension. The general conclusion summarizes the findings and future work. The appendices document policy training and the accepted WF-IoT paper.
\endgroup
